\documentclass[11pt]{article}

\usepackage[margin=1in]{geometry}
\usepackage{amsmath, amssymb}
\usepackage{xcolor}
\usepackage[skins,breakable]{tcolorbox}
\usepackage{hyperref}

\hypersetup{
    colorlinks=true,
    linkcolor=blue!60!black,
    urlcolor=blue!60!black,
    citecolor=blue!60!black
}

\newcommand{\SymPy}{SymPy}

% Public artifact repository; \bugbundle links a bug's bundle folder into it.
% TODO: set the public artifact repository URL before release.
\newcommand{\artifactrepo}{https://github.com/TODO-org/TODO-repo}
\newcommand{\bugbundle}[1]{\href{\artifactrepo/tree/main/bug-report/#1}{\nolinkurl{bug-report/#1/}}}

\newtcolorbox{counterexamplebox}{
  enhanced,
  colback=yellow!5,
  colframe=orange!70!black,
  arc=2mm,
  boxrule=0.8pt,
  left=4mm,
  right=4mm,
  top=3mm,
  bottom=3mm,
  before=\par\medskip\noindent,
  after=\par\medskip,
  before upper=\raggedright
}

\title{SymPy Bug Card}
\author{}
\date{}

\begin{document}

\maketitle

\section*{Bug 34: Branch Cut Error in Radical Solving}

\begin{counterexamplebox}
\textbf{Task.} Solve the principal-square-root equation over the complex domain:
\[
    \sqrt{1/x} = \frac{1}{\sqrt{x}}.
\]

\medskip
\textbf{SymPy's answer.}
\[
    \operatorname{solveset}\!\left(\sqrt{1/x}=1/\sqrt{x},\,x,\,\mathbb{C}\right)
    = \mathbb{C}\setminus\{0\}.
\]

\medskip
\textbf{Correct answer.}
\[
    \mathbb{C}\setminus(-\infty,0].
\]

\medskip
\textbf{Explanation.} Substituting \(x=-1\) gives
\(\sqrt{1/(-1)}=i\), but \(1/\sqrt{-1}=1/i=-i\). Thus the residual is \(2i\), so the negative real values included by \SymPy{} are not solutions.

\medskip
\textbf{Likely root cause.} The diagnosis traces the result through \texttt{\_solve\_radical} in \nolinkurl{sympy/solvers/solveset.py}: radical removal turns the checked equation into a broad infinite complement, and that complement is not rechecked against the original branch-sensitive equation. This is a new instance of a known failure mode.

\medskip
\textbf{Artifact (GitHub).} \bugbundle{bug-034-branch-cut-error-in-radical-solving}
\end{counterexamplebox}

\end{document}
